Computational thinking has emerged as a fundamental skill for the 21st century,
increasingly recognized as essential for all learners regardless of their future career paths~\cite{wing2006, grover2013}.
In Austria, this recognition materialized with the introduction of ``Digitale Grundbildung'' (Digital Basic Education)
as a mandatory subject for secondary education in 2022/23, reflecting a broader European trend
toward integrating computational concepts into primary and secondary curricula.
As a result, schools need programming tools that fit actual classroom conditions: they must be understandable for beginners,
manageable for teachers with limited programming experience, and suitable for an institutional context in which privacy, deployment simplicity, and curricular alignment all matter.


\section{Motivation and Problem Statement}
Many established block-based environments already support introductory programming quite well, yet they do not fully address this combination of requirements.
Scratch encourages creativity and experimentation, but its open-ended project structure makes it difficult to use for short, deterministically assessable classroom tasks.
Microsoft MakeCode offers a highly polished environment and a valuable bridge from blocks to text, but much of its design is centered on hardware integration and project-based exploration.
Code.org provides structured learning sequences and strong scaffolding, yet its deployment model and platform boundaries are less aligned with a privacy-first, self-hosted use case
in European school settings.
Teachers therefore face a practical gap between highly creative block environments on the one hand and the need for tightly guided, easy-to-manage learning tasks on the other.

This thesis addresses that gap through the design and implementation of BlockQuiz, a web-based platform for short block-programming exercises with automatic feedback. The central idea
is to combine the accessibility of block-based programming with a constrained exercise model that reduces cognitive load, supports clear learning objectives, and enables deterministic grading.
At the same time, the platform is designed for teachers who may not have a strong programming background and for schools that prefer local or self-hosted deployment over dependence on external services.


\section{Goals and Contributions}
The goal of the project is not to replace general-purpose creative environments, but to provide a complementary platform for guided practice. BlockQuiz focuses on short exercises in which
the available blocks, the expected interaction, and the grading criteria are defined in advance.
This structure allows the platform to deliver immediate and meaningful feedback while also giving teachers a manageable authoring workflow.
The thesis therefore contributes a constrained exercise model for beginner-oriented programming tasks, a browser-based sandboxing and grading architecture, a teacher-oriented content management approach
for exercise authoring, and a privacy-conscious deployment model that supports classroom use without relying on third-party cloud services.

\section{Target Audience}
The primary audience of the platform consists of children aged 8--12, a group for whom block-based interaction remains especially appropriate because it removes syntax barriers and keeps attention on
logic, sequencing, repetition, and conditional behavior. A second important audience consists of teachers who need to create, organize, and evaluate exercises without having to
program the platform itself.
Finally, the system also addresses administrators or school IT staff who need a solution that can be deployed and maintained with limited operational complexity.


\section{Scope and Limitations}
The implemented minimum viable product focuses on three exercise types: input/output tasks, turtle-graphics tasks, and grid-based robot tasks,
together with a basic course model, multilingual content, and a teacher-facing authoring workflow.
The platform emphasizes constrained tasks, client-side sandboxed execution, server-side authoritative grading for authenticated submissions, and Docker-based deployment.
It does not attempt to provide a complete learning management system, full classroom analytics, advanced institutional identity management, or a broad range of domain-specific simulators.
These omissions are intentional, as the thesis concentrates on the core software of a focused, privacy-conscious block-programming platform.

\section{Document Structure}
The remainder of this thesis follows the development of the platform.
Chapter \ref{chap:background} reviews the educational and technical background, discusses related platforms, and situates the work in existing literature on block-based programming and automatic
feedback.
Chapter \ref{chap:requirements} derives the requirements from the target group, the classroom context, and the intended deployment model.
Chapter \ref{chap:design} then explains the resulting system design and the main architectural decisions.
Chapter \ref{chap:implementation} describes the implementation of the platform in more detail, including the integration of Blockly, the execution sandbox,
and the grading framework.
Chapter \ref{chap:evaluation} presents a technical and heuristic evaluation of the prototype, including limitations and future validation needs.
Finally, Chapter \ref{chap:conclusion} summarizes the thesis, reflects on limitations, and identifies directions for future work.
